\section*{Bab \ref{ch:b2:muscle-movement} --- Kerutan Otot dan Fungsi Motorik}

\begin{solution}{exo:b2:muscle-movement:1}
Filamen tipisnya meluncur sepanjang filamen tebalnya, ditarik kepala
miosinnya; tidak ada filamen yang memendek. Pita I dan zona H-nya
menyempit dan cakram Z-nya mendekat; pita A-nya (yaitu filamen
tebalnya) menjaga panjangnya.
\end{solution}

\begin{solution}{exo:b2:muscle-movement:2}
(1) ATP terikat, kepalanya terlepas; ATP-nya terhidrolisis, kepalanya
tertegang. (2) Kepalanya mengikat aktin. (3) Fosfatnya dilepaskan,
langkah tenaganya, ADP-nya dilepaskan. (4) ATP yang baru mengikat dan
kepalanya terlepas. ATP diperlukan untuk \emph{melepaskan} kepalanya:
tanpanya tiap kepala tetap terikat dan ototnya terkunci --- rigor.
\end{solution}

\begin{solution}{exo:b2:muscle-movement:3}
Impuls sarafnya; pelepasan asetilkolin dan potensial lempeng-ujungnya
(\qty{0.6}{ms}); \omterm{prop:b2:neurons-synapses:ap}{potensial aksi} otot di seluruh serabutnya dan turun
lewat \omterm{def:b2:cell-differentiation:fibre}{tubulus T-nya} (\qty{1}{ms}); kalsium dilepaskan dari \omterm{def:b2:cell-differentiation:fibre}{retikulum
sarkoplasmanya} (\qty{1}{ms}); kalsiumnya mengikat troponin,
tropomiosinnya bergeser, kepalanya melekat (beberapa milidetik);
gayanya naik selama puluhan milidetik.
\end{solution}

\begin{solution}{exo:b2:muscle-movement:4}
\omterm{def:b2:muscle-movement:motorunit}{Unit motorik}: satu \omterm{def:b2:neurons-synapses:neuron}{neuron} motorik beserta seluruh serabut yang ia
sarafi. \omterm{def:b2:muscle-movement:motorunit}{Kedutan}: kerutan singkat dari sebuah impuls tunggal. \omterm{def:b2:muscle-movement:motorunit}{Tetanus}:
kerutan yang melebur dan berkelanjutan dari impuls yang tiba lebih
cepat daripada meluruhnya \omterm{def:b2:muscle-movement:motorunit}{kedutannya}. Gayanya dijenjangkan oleh laju
picu tiap unitnya dan oleh jumlah unit yang direkrut.
\end{solution}

\begin{solution}{exo:b2:muscle-movement:5}
Tiap separuh \omterm{def:b2:cell-differentiation:fibre}{sarkomernya} meluncur \qty{0.2}{\micro m} dalam
\qty{50}{ms}: \qty{4}{\micro m/s}; itu adalah 20 langkah \qty{10}{nm}
dalam \qty{50}{ms}, yaitu 400 tiap detik jika satu kepala terikat
sepanjang waktunya.
\end{solution}

\begin{solution}{exo:b2:muscle-movement:6}
\qty{2}{cm} dalam \qty{0.1}{s}: \qty{0.2}{m/s}, yaitu 2 panjang tiap
detik; tiap \omterm{def:b2:cell-differentiation:fibre}{sarkomernya} memendek $2/\num{40000}$ cm $= \qty{0.5}{\micro m}$ dalam
\qty{0.1}{s}, \qty{5}{\micro m/s}.
\end{solution}

\begin{solution}{exo:b2:muscle-movement:7}
$F/F_0 = 0.3125/(v/v_{\max} + 0.25) - 0.25$: pada 1, 3, dan 6 panjang
tiap detik ($0.1$, $0.3$, $0.6$ dari $v_{\max}$): $0.64$, $0.32$, $0.12$.
Dayanya $Fv$ (dalam $F_0\times$ panjang tiap detik): $0.64$, $0.95$, $0.71$
--- terbesar dekat 3 panjang tiap detik, yaitu kira-kira sepertiga $v_{\max}$.
\end{solution}

\begin{solution}{exo:b2:muscle-movement:8}
$80\times 30 = \qty{2400}{N}$; pada kakinya $2400\times 5/45 =
\qty{270}{N}$.
\end{solution}

\begin{solution}{exo:b2:muscle-movement:9}
Pada \qty{10}{Hz} impulsnya datang tiap \qty{100}{ms}, setelah kalsium
impuls sebelumnya lenyap: \omterm{def:b2:muscle-movement:motorunit}{kedutan} yang terpisah-pisah, yang
masing-masing bermula dari keadaan yang sebagian mengendur, sehingga
memberikan sebuah riak. Pada \qty{50}{Hz} impulsnya datang tiap
\qty{20}{ms}, sebelum kalsiumnya terpompa pergi: kalsiumnya tetap
tinggi dan gayanya melebur. \omterm{def:b2:muscle-movement:motorunit}{Tetanusnya} melampaui \omterm{def:b2:muscle-movement:motorunit}{kedutannya} karena
sebuah denyut kalsium tunggal terlalu singkat bagi jembatan silangnya
untuk meregangkan unsur elastik serinya dan membangkitkan gaya penuh;
kalsium yang berkelanjutan membiarkannya.
\end{solution}

\begin{solution}{exo:b2:muscle-movement:10}
Kerjanya \qty{20}{kJ}, energi ATP-nya \qty{80}{kJ}, yaitu $1.6$ mol
ATP. \qty{5}{s} pertamanya: $0.8$ mol \omterm{prop:b2:muscle-movement:energy}{kreatin fosfat}. Sisa $0.8$ mol
ATP-nya lewat glikolisis: $0.4$ mol glukosa, $0.8$ mol laktat --- di
dalam \qty{40}{L}, \qty{20}{mmol/L}.
\end{solution}

\begin{solution}{exo:b2:muscle-movement:11}
Kanal pelepas yang terhalang: tanpa kalsium, tanpa kerutan ---
kelumpuhan. Pompa yang terhalang: kalsiumnya tetap tinggi, ototnya
tidak dapat mengendur --- sebuah kontraktur. Troponin yang tak peka:
tropomiosinnya tidak pernah bergeser, tanpa kerutan. Tanpa kalsium di
dalam \omterm{def:b2:blood-circulation:blood}{darahnya}: otot rangkanya berkerut secara normal (kalsiumnya
berada di dalam); \omterm{def:b2:heart:anatomy}{jantungnya}, yang memerlukan kalsium luar untuk
memicu pelepasannya, berhenti.
\end{solution}

\begin{solution}{exo:b2:muscle-movement:12}
Kepalanya mengubah sampai separuh energi bebas sebuah ATP menjadi
kerja; sisanya adalah panas, dan rugi pompanya, rugi jembatan silang
yang melekat tanpa menarik pada kecepatan tinggi, serta rugi unsur
elastiknya membawa keseluruhannya menjadi seperempat --- yaitu daya
guna sebuah mesin bensin yang baik. Dua cara menjenjangkan gaya
membolehkan baik kendali yang halus pada usaha yang rendah (unit kecil
yang ditambahkan satu demi satu, masing-masing memicu lebih cepat atau
lebih lambat) maupun sebuah jangkauan yang luas (seribu kali lipat,
dari satu unit kecil pada laju \omterm{def:b2:muscle-movement:motorunit}{kedutan} sampai tiap unit dalam \omterm{def:b2:muscle-movement:motorunit}{tetanus});
sebuah sakelar tunggal akan memberikan satu gaya atau tidak sama sekali.
\end{solution}

\begin{solution}{pb:b2:muscle-movement:1}
\textbf{1.} $v = \sqrt{2\times 9.81\times 0.4} = \qty{2.8}{m/s}$;
$\tfrac12\times 70\times 2.8^{2} = \qty{275}{J}$.
\textbf{2.} Kerjanya $= 275 + 70\times 9.81\times 0.4 = \qty{550}{J}$ sepanjang
\qty{0.4}{m}: \qty{1370}{N}, yaitu dua kali beratnya.
\textbf{3.} $F_0 = 500\times 30 = \qty{15000}{N}$; pada kakinya,
\qty{1500}{N}.
\textbf{4.} \qty{4}{cm} dalam \qty{0.3}{s}: \qty{0.13}{m/s}, yaitu $1.1$ panjang tiap
detik, $0.14\,v_{\max}$.
\textbf{5.} $F/F_0 = 0.3125/(0.14 + 0.25) - 0.25 = 0.55$: \qty{8300}{N}
di dalam ototnya, \qty{830}{N} pada kakinya --- jauh kurang dari
\qty{1370}{N} yang dibutuhkan. Pemendekan ototnya saja tidak dapat
membuat lompatan ini. Gerakan lawannya memasok sisanya: ketika
pelompatnya merunduk, tendon yang teregang menyimpan energi elastik,
dan ototnya, yang berkerut nyaris secara isometrik pada gaya yang
tinggi, memuatinya; tendonnya lalu mengerut kembali lebih cepat
daripada yang dapat dicapai serabut mana pun ketika memendek.
\textbf{6.} $550/0.3 = \qty{1.8}{kW}$; \qty{92}{W/kg} otot pelurusnya.
\textbf{7.} $4/\num{48000}$ cm $= \qty{0.83}{\micro m}$ tiap \omterm{def:b2:cell-differentiation:fibre}{sarkomernya},
\qty{0.42}{\micro m} tiap separuhnya: 42 langkah \qty{10}{nm}.
\textbf{8.} \qty{0.42}{\micro m} dalam \qty{0.3}{s}: \qty{1.4}{\micro m/s};
140 langkah tiap detik jika terikat terus-menerus.
\textbf{9.} $500\times 6\times 10^{10}\times \num{48000}\times 300 =
4.3\times 10^{20}$ kepala.
\textbf{10.} $0.55\times \num{15000}\times 0.04 = \qty{330}{J}$; tiap
langkahnya $0.4\times 8\times 10^{-20} = \qty{3.2e-20}{J}$; $1.0\times
10^{22}$ langkah.
\textbf{11.} $1.0\times 10^{22}/4.3\times 10^{20} = 24$ langkah tiap
kepalanya, dari 42 yang tersedia: kira-kira \qty{60}{\%}. Cadangannya
kecil --- ototnya bekerja mendekati batasnya, seperti yang ditemukan
pertanyaan 5.
\textbf{12.} $1.0\times 10^{22}/6.0\times 10^{23} = 0.017$ mol,
\qty{8.7}{g}; energinya $0.017\times 50 = \qty{860}{J}$, yang
\qty{330}{J} darinya menjadi kerja: \qty{39}{\%}.
\textbf{13.} Lebih cepat daripada satu impuls tiap \qty{30}{ms}: di
atas \qty{33}{Hz}; kira-kira 10 impuls dalam \qty{0.3}{s}.
\textbf{14.} $9.9\times 10^{-6}\times 5\times 10^{-10}\times 6\times
10^{23} = 3\times 10^{9}$ ion bebas; $1.5\times 10^{9}$ ATP.
\textbf{15.} Serabutnya: $500/(\pi\times(3\times 10^{-3})^{2}) = 1.8\times
10^{7}$; ATP jembatan silang tiap serabut tiap impulsnya $1.0\times
10^{22}/(1.8\times 10^{7}\times 10) = 6\times 10^{13}$, yaitu empat puluh ribu
kali taksiran pompanya. Kalsium bebasnya adalah sisa yang kecil:
retikulumnya melepaskan sekitar \qty{0.1}{mmol/L}, yaitu sepuluh kali
kenaikan bebasnya, yang hampir semuanya seketika terikat troponin dan
penyangganya --- dan semuanya harus dipompa kembali.
\textbf{16.} Bagi sebuah lompatan maksimal, nyaris tiap unitnya dalam
beberapa puluh milidetik, yang kecil lebih dahulu dan yang besar
terakhir; sebuah langkah yang lembut merekrut barangkali sepersepuluh
sampai seperlimanya.
\textbf{17.} Unitnya direkrut menurut urutan ukurannya: pada usaha yang
rendah tiap unit baru menambahkan sebuah kenaikan yang kecil, pada
usaha yang tinggi unit yang tersisa berukuran besar dan masing-masing
menambahkan sebuah langkah yang besar.
\textbf{18.} Potensial lempeng-ujung tiap serabutnya menyetengah;
serabut yang potensialnya turun di bawah ambangnya tidak memicu,
sehingga \omterm{def:b2:muscle-movement:motorunit}{kedutan} dan \omterm{def:b2:muscle-movement:motorunit}{tetanusnya} lebih kecil, dan lompatannya kurang
jauh.
\textbf{19.} $5\times 20 = \qty{0.1}{mol}$ ATP: senilai enam lompatan.
\textbf{20.} $20\times 20 = \qty{0.4}{mol}$ \omterm{prop:b2:muscle-movement:energy}{kreatin fosfat}: kira-kira
23 lompatan.
\textbf{21.} Sepuluh lompatan memakai $0.17$ mol, yang belum
menghabiskan \omterm{prop:b2:muscle-movement:energy}{kreatin fosfatnya}; setelah dua puluhan lompatan,
glikolisisnya mengambil alih, dan laktat serta protonnya menumpuk.
\textbf{22.} \qty{1.8}{kW} kerja pada \qty{40}{\%} adalah \qty{4.6}{kW}
ATP, yaitu $0.09$ mol tiap detik, yang memerlukan $0.015$ mol oksigen
tiap detik --- \qty{0.34}{L/s}, \qty{21}{L/min}, yaitu sepuluh kali
yang dapat dibawa \omterm{def:b2:blood-circulation:blood}{darahnya}: lompatannya bersifat anaerob karena
keharusan, dan oksigennya datang sesudahnya.
\textbf{23.} $0.17$ mol ATP memerlukan $0.029$ mol oksigen,
\qty{0.64}{L}; jika dibayar selama semenit, \qty{640}{mL/min}, yaitu
dua setengah kali pemakaian istirahatnya.
\textbf{24.} ATP-nya difosforilasi ulang dari \omterm{prop:b2:muscle-movement:energy}{kreatin fosfatnya} dalam
hitungan milidetik, sehingga kepekatannya nyaris tidak turun; rigor
memerlukan ATP yang mendekati nol, yang tidak pernah dicapai sebuah
otot hidup yang memiliki cadangan apa pun.
\textbf{25.} 42 langkah yang tersedia tiap separuh \omterm{def:b2:cell-differentiation:fibre}{sarkomernya}, 24 yang
dibutuhkan; $F/F_0 = 0.55$ pada kecepatan lompatannya; $0.017$ mol,
\qty{8.7}{g}, ATP.
\end{solution}
